% from ega4-2-fr.tex
\oldpage[IV-2]{134}

on ait $\operatorname{Tor}^A_i(M,N)=0$ (0, 17.3.1 et 17.2.6). Mais on a le
lemme suivant (où les $A_\alpha$ sont de nouveau quelconques):

\begin{lemma}[5.13.7.1]
\label{IV.2.s5.13.7.1-fr}
Pour tout $A$-module $M$ de présentation finie, il existe un indice $\alpha$
et un $A_\alpha$-module de présentation finie $M_\alpha$ tel que $M$ soit
isomorphe à $M_\alpha\otimes_{A_\alpha}A$.
\end{lemma}

En effet, on peut supposer que $M=A^r/P$, où $P$ est un sous-$A$-module de
type fini de $A^r$. Comme $A^r=\varinjlim(A_\alpha^r)$ et que $P$ est de type
fini, il existe un $\alpha$ et un sous-$A_\alpha$-module de type fini
$P_\alpha$ de $A_\alpha^r$ tel que $P$ soit l'image canonique de
$P_\alpha\otimes_{A_\alpha}A$; l'exactitude à droite du produit tensoriel
montre donc que $M$ est isomorphe à
$(A_\alpha^r/P_\alpha)\otimes_{A_\alpha}A$.

Ce lemme étant établi, $M$ et $N$ sont de présentation finie, puisque $A$ est
noethérien; écrivons donc
$M=M_\alpha\otimes_{A_\alpha}A$,
$N=N_\alpha\otimes_{A_\alpha}A$; en vertu de l'hypothèse, $A$ est un
$A_\alpha$-module plat $(0_{\mathrm I}, 6.1.2)$, donc on a
\[
\operatorname{Tor}^A_i(M,N)
=\operatorname{Tor}^{A_\alpha}_i(M_\alpha,N_\alpha)
 \otimes_{A_\alpha}A
\]
à un isomorphisme près (III, 6.3.9); comme par hypothèse $A_\alpha$ est
régulier, il existe $i_0$ tel que
$\operatorname{Tor}^{A_\alpha}_i(M_\alpha,N_\alpha)=0$ pour $i\geq i_0$.
C.Q.F.D.

\begin{corollary}[5.13.8]
\label{IV.2.s5.13.8-fr}
Supposons que les $A_\alpha$ soient noethériens et vérifient la propriété
$(R_k)$ pour un entier $k\geq0$. Supposons en outre que pour
$\alpha\leq\beta$, $\varphi_{\beta\alpha}$ fasse de $A_\beta$ un
$A_\alpha$-module plat, que $A=\varinjlim A_\alpha$ soit noethérien et que,
pour tout $\alpha$, ${}^a\varphi_\alpha$ soit un homéomorphisme de
$\operatorname{Spec}(A)$ sur $\operatorname{Spec}(A_\alpha)$. Alors $A$
vérifie $(R_k)$.
\end{corollary}

Soit en effet $\mathfrak p$ un idéal premier de $A$ et, pour tout $\alpha$,
posons $\mathfrak p_\alpha=\varphi_\alpha^{-1}(\mathfrak p)$, de sorte que
$\mathfrak p=\varinjlim\mathfrak p_\alpha$ et
$A_{\mathfrak p}=\varinjlim(A_\alpha)_{\mathfrak p_\alpha}$ (5.13.3). Il
résulte de (I, 2.4.2) et de l'hypothèse que le morphisme
\[
\operatorname{Spec}(A_{\mathfrak p})
\longrightarrow\operatorname{Spec}((A_\alpha)_{\mathfrak p_\alpha})
\]
est un homéomorphisme surjectif pour tout $\alpha$; on a donc
$\dim(A_{\mathfrak p})=\dim((A_\alpha)_{\mathfrak p_\alpha})$ pour tout
$\alpha$. Supposons alors que $\dim(A_{\mathfrak p})\leq k$; on a donc
$\dim((A_\alpha)_{\mathfrak p_\alpha})\leq k$, et l'hypothèse entraîne que
$(A_\alpha)_{\mathfrak p_\alpha}$ est un anneau régulier; comme en outre, pour
$\alpha\leq\beta$, $(A_\beta)_{\mathfrak p_\beta}$ est un
$(A_\alpha)_{\mathfrak p_\alpha}$-module plat $(0_{\mathrm I}, 6.3.2)$, on
conclut de (5.13.7) que $A_{\mathfrak p}$ est régulier, ce qui démontre le
corollaire.

\section{Morphismes plats de préschémas localement noethériens}

Soient $X$, $Y$ deux préschémas localement noethériens, $f:X\to Y$ un
morphisme. Pour tout $y\in Y$, la fibre
$f^{-1}(y)=X\times_Y\operatorname{Spec}(k(y))$ est aussi un préschéma
localement noethérien, car il suffit de le voir lorsque
$Y=\operatorname{Spec}(A)$ et $X=\operatorname{Spec}(B)$ sont des spectres
d'anneaux noethériens, et alors $B\otimes_A k(y)$ est anneau de fractions de
l'anneau quotient $B/\mathfrak j_yB$, donc est noethérien. Nous nous proposons
en premier lieu dans ce paragraphe de mettre en relation les propriétés de
$X$, de $Y$ et des fibres $f^{-1}(y)$ (où $y$ parcourt $Y$), sous l'hypothèse
que le morphisme $f$ est plat. Les questions traitées se ramèneront à l'étude
des relations entre les anneaux locaux noethériens
$\mathcal O_y$, $\mathcal O_x$ et
$\mathcal O_x\otimes_{\mathcal O_y}k(y)$, l'homomorphisme
$\mathcal O_y\to\mathcal O_x$ étant local et faisant de $\mathcal O_x$ un
$\mathcal O_y$-module plat. Dans les n\textsuperscript{os} 6.11 à 6.13 (assez
séparés du reste du paragraphe, de par leur nature « absolue » plutôt que
« relative »), nous appliquons certains des résultats précédents pour trouver
des

